% ======================================================================
\section{Podsumowanie}
% ======================================================================

\begin{frame}{Kluczowe idee}
  \begin{enumerate}
    \item<1-> \textbf{Rekurencja à la Strassen} z tożsamością Schönhagego (10 mnożeń zamiast 13):
      $L$ poziomów daje naraz $\binom Lm$ cienkich iloczynów $N_0\times D\times N_0$.
    \item<2-> \textbf{Błąd $E$ ginie}, bo jego jednomiany mieszają zmienne wewnętrzne z zewnętrznymi (Lemat 9).
    \item<3-> \textbf{Kafelkowanie siatką}: każdy wpis $XY$ to jedno wyjście jednego $\Full$, a tablica wejściowa zależy tylko od pasa.
    \item<4-> \textbf{Kodowanie raz na pas}: kodowanie zależy tylko od $X$ albo tylko od $Y$, a $N\ge D^{18}$ czyni je pomijalnym.
    \item<5-> \textbf{Przycinanie}: odwiedzamy tylko $\Leaves(U)$. Każde wywołanie raz, $O(L^2)$ na liść (Lemat 10).
    \item<6-> \textbf{Liczenie liści według rzędu}: prywatne liście ($d=0$) plus liście współdzielone.
      Małe rzędy liczymy na napis, a duże globalnie ($\beta_d$ maleje geometrycznie przy $L=19m$).
      Stąd $O(M/D^{1/18})$ liści na kafelek przy gęstości $1/\sqrt D$.
  \end{enumerate}
\end{frame}

\begin{frame}[c]
  \centering
  \definecolor{claude}{HTML}{D97757}
  \definecolor{diabel}{HTML}{8E44AD}
  \begin{tikzpicture}[scale=1.25]
    % rogi jak w emoji z diabłem (rysowane pod promieniami)
    \foreach \s in {1,-1}
      \fill[diabel,xscale=\s] (0.16,0.70) .. controls (0.45,0.92) and (0.72,1.12) .. (0.84,1.50)
        .. controls (1.02,1.18) and (0.96,0.80) .. (0.60,0.50) -- cycle;
    \foreach \kat/\dl in {74/1.00,116/1.02,146/0.94,180/0.95,212/0.90,239/0.92,
                          272/0.95,311/1.00,343/0.88,5/0.88,36/0.96,95/0.80}
      \fill[claude] ($(0,0)+(\kat+90:0.035)$) -- ($(\kat:\dl)+(\kat+90:0.075)$)
        arc[start angle=\kat+90,end angle=\kat-90,radius=0.075]
        -- ($(0,0)+(\kat-90:0.035)$) -- cycle;
    \fill[claude] (0,0) circle (0.12);
  \end{tikzpicture}

  \vspace{0.8em}
  {\Huge\textcolor{structure}{\textbf{Koniec}}}

  \vspace{1em}
  {\large Stanisław Macura \quad Michał Chłopicki}
\end{frame}
